\documentclass[12pt,a4paper]{article}
\usepackage[utf8]{inputenc}
\usepackage[T1]{fontenc}
\usepackage[french]{babel}
\usepackage{geometry}
\usepackage{graphicx}
\usepackage{amsmath}
\usepackage{amssymb}
\usepackage{booktabs}
\usepackage{longtable}
\usepackage{listings}
\usepackage{xcolor}
\usepackage{hyperref}
\usepackage{tikz}
\usetikzlibrary{er,positioning,shapes}

\geometry{margin=2.5cm}

\lstset{
    basicstyle=\ttfamily\small,
    keywordstyle=\color{blue}\bfseries,
    commentstyle=\color{gray},
    stringstyle=\color{red},
    breaklines=true,
    frame=single,
    backgroundcolor=\color{gray!10}
}

\title{
    \textbf{PROJET DE MODULE}\\
    \vspace{0.5cm}
    \Large Smart City Analytics Platform\\
    \large Plateforme Intelligente de Gestion des Données Urbaines\\
    \vspace{1cm}
    \normalsize Neo-Sousse 2030
}

\author{
    Section : IA2\\
    \vspace{0.5cm}
    Université de Sousse\\
    École Nationale d'Ingénieurs de Sousse\\
    Département d'Informatique Appliquée
}

\date{Année Universitaire 2025/2026}

\begin{document}

\maketitle
\thispagestyle{empty}
\newpage

\tableofcontents
\newpage

%==============================================================================
\section{Introduction}
%==============================================================================

Ce rapport présente la conception et l'implémentation d'une plateforme de gestion des données urbaines pour la métropole \textbf{Neo-Sousse 2030}. L'objectif est de collecter, analyser et visualiser les données IoT pour améliorer la qualité de vie des citoyens et réduire l'empreinte écologique.

\subsection{Contexte et Enjeux}
Face à la croissance démographique et aux défis environnementaux, la métropole Neo-Sousse s'est engagée dans une transformation digitale. Cette plateforme permet de :
\begin{itemize}
    \item Gérer les capteurs IoT (qualité de l'air, trafic, énergie, déchets, éclairage)
    \item Suivre les interventions de maintenance
    \item Analyser l'engagement des citoyens
    \item Optimiser les trajets des véhicules autonomes
\end{itemize}

\subsection{Stack Technique}
\begin{itemize}
    \item \textbf{SGBD} : MySQL (relationnel, conformité ACID)
    \item \textbf{Backend} : Python Django + Django REST Framework
    \item \textbf{Frontend} : React.js (interface basique)
    \item \textbf{ETL} : Scripts Python intégrés
\end{itemize}

%==============================================================================
\section{Partie A : Modélisation Théorique}
%==============================================================================

%------------------------------------------------------------------------------
\subsection{Modélisation Conceptuelle - Diagramme E/A}
%------------------------------------------------------------------------------

Le diagramme Entité-Association ci-dessous représente l'ensemble des besoins fonctionnels du système.

\subsubsection{Entités Identifiées}

\begin{table}[h]
\centering
\begin{tabular}{|l|l|p{8cm}|}
\hline
\textbf{Entité} & \textbf{Identifiant} & \textbf{Description} \\
\hline
ARRONDISSEMENT & id & Zones géographiques de la métropole \\
\hline
PROPRIETAIRE & id & Propriétaires des capteurs (municipalité ou privé) \\
\hline
CAPTEUR & uuid & Capteurs IoT déployés dans la ville \\
\hline
MESURE & id & Relevés collectés par les capteurs \\
\hline
TECHNICIEN & id & Personnel de maintenance \\
\hline
INTERVENTION & id & Opérations de maintenance sur les capteurs \\
\hline
CITOYEN & id & Citoyens engagés dans la démarche écologique \\
\hline
CONSULTATION & id & Historique des consultations de données \\
\hline
VEHICULE & id & Véhicules autonomes municipaux \\
\hline
TRAJET & id & Trajets effectués par les véhicules \\
\hline
\end{tabular}
\caption{Liste des entités du système}
\end{table}

\subsubsection{Diagramme E/A}

\begin{center}
\begin{tikzpicture}[
    entity/.style={rectangle, draw, fill=blue!20, minimum width=2.5cm, minimum height=0.8cm},
    relationship/.style={diamond, draw, fill=green!20, aspect=2},
    attribute/.style={ellipse, draw, fill=yellow!20, font=\small}
]

% Entités principales
\node[entity] (arr) at (0,8) {ARRONDISSEMENT};
\node[entity] (prop) at (6,8) {PROPRIETAIRE};
\node[entity] (cap) at (3,5) {CAPTEUR};
\node[entity] (mes) at (0,2) {MESURE};
\node[entity] (tech) at (9,5) {TECHNICIEN};
\node[entity] (inter) at (6,2) {INTERVENTION};
\node[entity] (cit) at (12,8) {CITOYEN};
\node[entity] (cons) at (12,5) {CONSULTATION};
\node[entity] (veh) at (0,-1) {VEHICULE};
\node[entity] (traj) at (4,-1) {TRAJET};

% Relations
\node[relationship] (r1) at (1.5,6.5) {localise};
\node[relationship] (r2) at (4.5,6.5) {possède};
\node[relationship] (r3) at (1.5,3.5) {génère};
\node[relationship] (r4) at (4.5,3.5) {concerne};
\node[relationship] (r5) at (7.5,3.5) {intervient};
\node[relationship] (r6) at (8,3) {valide};
\node[relationship] (r7) at (12,6.5) {effectue};
\node[relationship] (r8) at (2,-1) {réalise};

% Connexions
\draw (arr) -- (r1) -- (cap);
\draw (prop) -- (r2) -- (cap);
\draw (cap) -- (r3) -- (mes);
\draw (cap) -- (r4) -- (inter);
\draw (tech) -- (r5) -- (inter);
\draw (tech) -- (r6) -- (inter);
\draw (cit) -- (r7) -- (cons);
\draw (veh) -- (r8) -- (traj);

\end{tikzpicture}
\end{center}

%------------------------------------------------------------------------------
\subsection{Cardinalités et Contraintes}
%------------------------------------------------------------------------------

\subsubsection{Cardinalités des Associations}

\begin{table}[h]
\centering
\begin{tabular}{|l|c|c|l|}
\hline
\textbf{Association} & \textbf{Entité 1} & \textbf{Entité 2} & \textbf{Justification} \\
\hline
ARRONDISSEMENT-CAPTEUR & (1,1) & (0,n) & Un capteur est dans un seul arrondissement \\
\hline
PROPRIETAIRE-CAPTEUR & (1,1) & (0,n) & Un capteur a un seul propriétaire \\
\hline
CAPTEUR-MESURE & (1,1) & (0,n) & Une mesure provient d'un seul capteur \\
\hline
CAPTEUR-INTERVENTION & (1,1) & (0,n) & Une intervention concerne un capteur \\
\hline
TECHNICIEN-INTERVENTION & (1,1) & (0,n) & Intervenant obligatoire \\
\hline
TECHNICIEN-INTERVENTION & (1,1) & (0,n) & Validateur obligatoire \\
\hline
CITOYEN-CONSULTATION & (1,1) & (0,n) & Une consultation est faite par un citoyen \\
\hline
VEHICULE-TRAJET & (1,1) & (0,n) & Un trajet est effectué par un véhicule \\
\hline
\end{tabular}
\caption{Cardinalités des associations}
\end{table}

\subsubsection{Contrainte Métier Critique}

\begin{center}
\fbox{
\parbox{0.8\textwidth}{
\textbf{Règle Métier :} Chaque intervention nécessite \textbf{au moins deux techniciens différents} :
\begin{itemize}
    \item Un \textbf{technicien intervenant} qui effectue l'opération
    \item Un \textbf{technicien validateur} (humain ou système IA) qui valide
\end{itemize}
\textbf{Contrainte :} $technicien\_intervenant\_id \neq technicien\_validateur\_id$
}}
\end{center}

%------------------------------------------------------------------------------
\subsection{Modèle Relationnel Normalisé}
%------------------------------------------------------------------------------

\subsubsection{Schéma Relationnel}

\textbf{ARRONDISSEMENT} (\underline{id}, nom, code\_postal, superficie\_km2, population)

\textbf{PROPRIETAIRE} (\underline{id}, nom, type, contact\_email, contact\_telephone, adresse)

\textbf{CAPTEUR} (\underline{uuid}, type, latitude, longitude, statut, date\_installation, description, \#proprietaire\_id, \#arrondissement\_id)

\textbf{MESURE} (\underline{id}, \#capteur\_uuid, timestamp, valeurs, indice\_pollution, qualite\_air)

\textbf{TECHNICIEN} (\underline{id}, matricule, nom, prenom, email, telephone, specialite, type\_validation, actif)

\textbf{INTERVENTION} (\underline{id}, \#capteur\_uuid, date\_heure, nature, description, duree\_minutes, cout, reduction\_co2, statut, \#technicien\_intervenant\_id, \#technicien\_validateur\_id)

\textbf{CITOYEN} (\underline{id}, identifiant\_unique, nom, prenom, email, telephone, date\_naissance, adresse, score\_engagement, preference\_mobilite, notifications\_actives)

\textbf{CONSULTATION} (\underline{id}, \#citoyen\_id, timestamp, type\_donnee, zone\_consultee, duree\_consultation\_secondes, source\_consultation)

\textbf{VEHICULE} (\underline{id}, plaque, type, energie, marque, modele, annee\_mise\_service, capacite\_passagers, autonomie\_km, statut)

\textbf{TRAJET} (\underline{id}, \#vehicule\_id, origine, destination, depart, arrivee, distance\_km, economie\_co2, nombre\_passagers, statut)

\vspace{0.5cm}
\textit{Légende : \underline{clé primaire}, \#clé étrangère}

%------------------------------------------------------------------------------
\subsection{Analyse des Redondances}
%------------------------------------------------------------------------------

\subsubsection{Redondances Potentielles Identifiées}

\begin{enumerate}
    \item \textbf{Nom de l'arrondissement dans CAPTEUR} : Évitée en utilisant une clé étrangère vers ARRONDISSEMENT au lieu de stocker le nom directement.
    
    \item \textbf{Coordonnées du propriétaire dans CAPTEUR} : Évitée par la relation avec PROPRIETAIRE.
    
    \item \textbf{Informations technicien dans INTERVENTION} : Évitée par les clés étrangères vers TECHNICIEN.
\end{enumerate}

\subsubsection{Solutions Appliquées}

Toutes les redondances ont été éliminées par :
\begin{itemize}
    \item Création de tables de référence (ARRONDISSEMENT, PROPRIETAIRE, TECHNICIEN)
    \item Utilisation de clés étrangères pour les associations
    \item Séparation des données temporelles (MESURE, TRAJET) des données de référence
\end{itemize}

%------------------------------------------------------------------------------
\subsection{Dépendances Fonctionnelles}
%------------------------------------------------------------------------------

\subsubsection{Graphe Minimal des Dépendances Fonctionnelles}

\textbf{Table CAPTEUR :}
\begin{align*}
uuid &\rightarrow type, latitude, longitude, statut, date\_installation, description \\
uuid &\rightarrow proprietaire\_id, arrondissement\_id
\end{align*}

\textbf{Table INTERVENTION :}
\begin{align*}
id &\rightarrow capteur\_uuid, date\_heure, nature, description \\
id &\rightarrow duree\_minutes, cout, reduction\_co2, statut \\
id &\rightarrow technicien\_intervenant\_id, technicien\_validateur\_id
\end{align*}

\textbf{Table MESURE :}
\begin{align*}
id &\rightarrow capteur\_uuid, timestamp, valeurs, indice\_pollution, qualite\_air
\end{align*}

\textbf{Table CITOYEN :}
\begin{align*}
id &\rightarrow identifiant\_unique, nom, prenom, email, ... \\
identifiant\_unique &\rightarrow id \quad \text{(dépendance réciproque)}\\
email &\rightarrow id \quad \text{(email unique)}
\end{align*}

\textbf{Table TRAJET :}
\begin{align*}
id &\rightarrow vehicule\_id, origine, destination, depart, arrivee \\
id &\rightarrow distance\_km, economie\_co2, nombre\_passagers, statut
\end{align*}

%------------------------------------------------------------------------------
\subsection{Clés et Identifiants}
%------------------------------------------------------------------------------

\begin{table}[h]
\centering
\begin{tabular}{|l|l|l|}
\hline
\textbf{Table} & \textbf{Clé Primaire} & \textbf{Clés Candidates / Uniques} \\
\hline
ARRONDISSEMENT & id & code\_postal (UNIQUE) \\
\hline
PROPRIETAIRE & id & - \\
\hline
CAPTEUR & uuid & - \\
\hline
MESURE & id & - \\
\hline
TECHNICIEN & id & matricule (UNIQUE) \\
\hline
INTERVENTION & id & - \\
\hline
CITOYEN & id & identifiant\_unique (UNIQUE), email (UNIQUE) \\
\hline
CONSULTATION & id & - \\
\hline
VEHICULE & id & plaque (UNIQUE) \\
\hline
TRAJET & id & - \\
\hline
\end{tabular}
\caption{Clés primaires et candidates}
\end{table}

\subsubsection{Clés Étrangères}

\begin{table}[h]
\centering
\begin{tabular}{|l|l|l|}
\hline
\textbf{Table} & \textbf{Clé Étrangère} & \textbf{Référence} \\
\hline
CAPTEUR & proprietaire\_id & PROPRIETAIRE(id) \\
\hline
CAPTEUR & arrondissement\_id & ARRONDISSEMENT(id) \\
\hline
MESURE & capteur\_uuid & CAPTEUR(uuid) \\
\hline
INTERVENTION & capteur\_uuid & CAPTEUR(uuid) \\
\hline
INTERVENTION & technicien\_intervenant\_id & TECHNICIEN(id) \\
\hline
INTERVENTION & technicien\_validateur\_id & TECHNICIEN(id) \\
\hline
CONSULTATION & citoyen\_id & CITOYEN(id) \\
\hline
TRAJET & vehicule\_id & VEHICULE(id) \\
\hline
\end{tabular}
\caption{Clés étrangères du système}
\end{table}

%------------------------------------------------------------------------------
\subsection{Formes Normales}
%------------------------------------------------------------------------------

\subsubsection{Vérification de la 1ère Forme Normale (1FN)}

\textbf{Critères :}
\begin{itemize}
    \item[$\checkmark$] Toutes les colonnes contiennent des valeurs atomiques
    \item[$\checkmark$] Pas de groupes répétitifs
    \item[$\checkmark$] Chaque table possède une clé primaire
\end{itemize}

\textbf{Exemple :} Les coordonnées GPS sont séparées en deux colonnes (latitude, longitude) au lieu d'un seul champ "position".

\subsubsection{Vérification de la 2ème Forme Normale (2FN)}

\textbf{Critères :}
\begin{itemize}
    \item[$\checkmark$] Respect de la 1FN
    \item[$\checkmark$] Tous les attributs non-clés dépendent entièrement de la clé primaire
\end{itemize}

\textbf{Analyse :} Aucune table n'a de clé primaire composée, donc la 2FN est automatiquement respectée.

\subsubsection{Vérification de la 3ème Forme Normale (3FN)}

\textbf{Critères :}
\begin{itemize}
    \item[$\checkmark$] Respect de la 2FN
    \item[$\checkmark$] Aucune dépendance transitive
\end{itemize}

\textbf{Exemples de conformité :}
\begin{itemize}
    \item Le nom de l'arrondissement n'est pas stocké dans CAPTEUR (évite nom $\rightarrow$ code\_postal)
    \item Les coordonnées du propriétaire ne sont pas dans CAPTEUR
    \item Le nom du technicien n'est pas répété dans INTERVENTION
\end{itemize}

\begin{center}
\fbox{\textbf{Conclusion : Toutes les tables sont en 3FN}}
\end{center}

%------------------------------------------------------------------------------
\subsection{Décomposition BCNF (Boyce-Codd Normal Form)}
%------------------------------------------------------------------------------

\subsubsection{Analyse BCNF}

Une table est en BCNF si pour toute dépendance fonctionnelle $X \rightarrow Y$, X est une super-clé.

\textbf{Vérification table CITOYEN :}
\begin{itemize}
    \item $id \rightarrow$ tous les attributs $\checkmark$ (id est clé primaire)
    \item $identifiant\_unique \rightarrow id$ $\checkmark$ (identifiant\_unique est clé candidate)
    \item $email \rightarrow id$ $\checkmark$ (email est clé candidate)
\end{itemize}

\textbf{Vérification table INTERVENTION :}
\begin{itemize}
    \item $id \rightarrow$ tous les attributs $\checkmark$ (id est clé primaire)
    \item Pas de dépendance partielle
\end{itemize}

\begin{center}
\fbox{\textbf{Conclusion : Toutes les tables respectent la BCNF}}
\end{center}

Aucune décomposition supplémentaire n'est nécessaire car le schéma initial a été conçu en respectant directement la BCNF.

%==============================================================================
\section{Partie B : Implémentation Pratique}
%==============================================================================

%------------------------------------------------------------------------------
\subsection{Scripts SQL de Génération des Tables}
%------------------------------------------------------------------------------

\subsubsection{Création de la Base de Données}

\begin{lstlisting}[language=SQL]
CREATE DATABASE IF NOT EXISTS smart_city_db
    CHARACTER SET utf8mb4
    COLLATE utf8mb4_unicode_ci;

USE smart_city_db;
\end{lstlisting}

\subsubsection{Table ARRONDISSEMENT}

\begin{lstlisting}[language=SQL]
CREATE TABLE arrondissement (
    id INT AUTO_INCREMENT PRIMARY KEY,
    nom VARCHAR(100) NOT NULL,
    code_postal VARCHAR(10) NOT NULL UNIQUE,
    superficie_km2 DECIMAL(10, 2),
    population INT,
    created_at TIMESTAMP DEFAULT CURRENT_TIMESTAMP
);
\end{lstlisting}

\subsubsection{Table CAPTEUR}

\begin{lstlisting}[language=SQL]
CREATE TABLE capteur (
    uuid CHAR(36) PRIMARY KEY,
    type ENUM('AIR', 'TRAFIC', 'ENERGIE', 'DECHETS', 'ECLAIRAGE') NOT NULL,
    latitude DECIMAL(10, 8) NOT NULL,
    longitude DECIMAL(11, 8) NOT NULL,
    statut ENUM('ACTIF', 'MAINTENANCE', 'HORS_SERVICE') DEFAULT 'ACTIF',
    date_installation DATE NOT NULL,
    proprietaire_id INT NOT NULL,
    arrondissement_id INT NOT NULL,
    
    FOREIGN KEY (proprietaire_id) REFERENCES proprietaire(id),
    FOREIGN KEY (arrondissement_id) REFERENCES arrondissement(id)
);
\end{lstlisting}

\subsubsection{Table INTERVENTION avec Contrainte 2 Techniciens}

\begin{lstlisting}[language=SQL]
CREATE TABLE intervention (
    id INT AUTO_INCREMENT PRIMARY KEY,
    capteur_uuid CHAR(36) NOT NULL,
    date_heure DATETIME NOT NULL,
    nature ENUM('PREDICTIVE', 'CORRECTIVE', 'CURATIVE') NOT NULL,
    duree_minutes INT NOT NULL,
    cout DECIMAL(10, 2) DEFAULT 0,
    reduction_co2 DECIMAL(10, 4),
    statut ENUM('PLANIFIEE', 'EN_COURS', 'TERMINEE', 'ANNULEE'),
    
    -- Contrainte: 2 techniciens obligatoires
    technicien_intervenant_id INT NOT NULL,
    technicien_validateur_id INT NOT NULL,
    
    FOREIGN KEY (capteur_uuid) REFERENCES capteur(uuid),
    FOREIGN KEY (technicien_intervenant_id) REFERENCES technicien(id),
    FOREIGN KEY (technicien_validateur_id) REFERENCES technicien(id),
    
    -- Contrainte: techniciens differents
    CONSTRAINT chk_techniciens_differents 
        CHECK (technicien_intervenant_id != technicien_validateur_id)
);
\end{lstlisting}

%------------------------------------------------------------------------------
\subsection{Vues et Requêtes Métiers}
%------------------------------------------------------------------------------

\subsubsection{Question 1 : Zones les plus polluées (24h)}

\begin{lstlisting}[language=SQL]
SELECT 
    a.nom AS arrondissement,
    AVG(m.indice_pollution) AS pollution_moyenne,
    COUNT(m.id) AS nombre_mesures
FROM arrondissement a
JOIN capteur c ON a.id = c.arrondissement_id
JOIN mesure m ON c.uuid = m.capteur_uuid
WHERE c.type = 'AIR'
  AND m.timestamp >= NOW() - INTERVAL 24 HOUR
GROUP BY a.id
ORDER BY pollution_moyenne DESC;
\end{lstlisting}

\subsubsection{Question 2 : Taux de disponibilité par arrondissement}

\begin{lstlisting}[language=SQL]
SELECT 
    a.nom AS arrondissement,
    COUNT(*) AS total_capteurs,
    SUM(CASE WHEN c.statut = 'ACTIF' THEN 1 ELSE 0 END) AS actifs,
    ROUND(SUM(CASE WHEN c.statut = 'ACTIF' THEN 1 ELSE 0 END) 
          * 100.0 / COUNT(*), 2) AS taux_disponibilite
FROM arrondissement a
JOIN capteur c ON a.id = c.arrondissement_id
GROUP BY a.id;
\end{lstlisting}

\subsubsection{Question 3 : Citoyens les plus engagés}

\begin{lstlisting}[language=SQL]
SELECT 
    c.identifiant_unique,
    CONCAT(c.prenom, ' ', c.nom) AS nom_complet,
    c.score_engagement,
    c.preference_mobilite,
    COUNT(co.id) AS nombre_consultations
FROM citoyen c
LEFT JOIN consultation co ON c.id = co.citoyen_id
GROUP BY c.id
ORDER BY c.score_engagement DESC
LIMIT 10;
\end{lstlisting}

\subsubsection{Question 4 : Interventions prédictives ce mois}

\begin{lstlisting}[language=SQL]
SELECT 
    COUNT(*) AS nombre_interventions,
    SUM(reduction_co2) AS economie_co2_totale,
    SUM(cout) AS cout_total
FROM intervention
WHERE nature = 'PREDICTIVE'
  AND statut = 'TERMINEE'
  AND MONTH(date_heure) = MONTH(CURRENT_DATE)
  AND YEAR(date_heure) = YEAR(CURRENT_DATE);
\end{lstlisting}

\subsubsection{Question 5 : Trajets avec meilleure réduction CO2}

\begin{lstlisting}[language=SQL]
SELECT 
    v.plaque,
    v.type AS type_vehicule,
    v.energie,
    t.origine,
    t.destination,
    t.distance_km,
    t.economie_co2
FROM trajet t
JOIN vehicule v ON t.vehicule_id = v.id
WHERE t.statut = 'TERMINE'
ORDER BY t.economie_co2 DESC
LIMIT 20;
\end{lstlisting}

%------------------------------------------------------------------------------
\subsection{Jeux de Données}
%------------------------------------------------------------------------------

\begin{table}[h]
\centering
\begin{tabular}{|l|c|}
\hline
\textbf{Entité} & \textbf{Nombre d'enregistrements} \\
\hline
Arrondissements & 5 \\
Propriétaires & 4 \\
Techniciens & 10 \\
Capteurs & 50 \\
Citoyens & 30 \\
Véhicules & 20 \\
Mesures & 500 \\
Interventions & 100 \\
Trajets & 200 \\
Consultations & 325 \\
\hline
\end{tabular}
\caption{Volume des données de test}
\end{table}

%------------------------------------------------------------------------------
\subsection{API REST (Backend Django)}
%------------------------------------------------------------------------------

\begin{table}[h]
\centering
\begin{tabular}{|l|l|l|}
\hline
\textbf{Méthode} & \textbf{Endpoint} & \textbf{Description} \\
\hline
GET & /api/capteurs/ & Liste des capteurs \\
GET & /api/capteurs/zones\_polluees\_24h/ & Zones polluées \\
GET & /api/capteurs/disponibilite\_par\_arrondissement/ & Taux disponibilité \\
GET & /api/interventions/ & Liste des interventions \\
GET & /api/citoyens/top\_engages/ & Top citoyens \\
GET & /api/trajets/top\_trajets\_co2/ & Meilleurs trajets CO2 \\
\hline
\end{tabular}
\caption{Endpoints API principaux}
\end{table}

%------------------------------------------------------------------------------
\subsection{Scripts ETL}
%------------------------------------------------------------------------------

\begin{lstlisting}[language=Python]
# Commande Django pour generer les donnees
python manage.py seed_data    # Donnees initiales
python manage.py run_etl      # Mesures, interventions, trajets
\end{lstlisting}

%==============================================================================
\section{Conclusion}
%==============================================================================

Ce projet a permis de concevoir et implémenter une plateforme complète de gestion des données urbaines respectant :

\begin{enumerate}
    \item Une \textbf{modélisation conceptuelle} rigoureuse (diagramme E/A)
    \item Un \textbf{schéma relationnel normalisé} en 3FN/BCNF
    \item Des \textbf{contraintes d'intégrité} fortes (FK, CHECK, UNIQUE)
    \item La \textbf{contrainte métier} des 2 techniciens par intervention
    \item Des \textbf{requêtes analytiques} répondant aux questions métiers
    \item Une \textbf{API REST} complète avec Django
\end{enumerate}

\end{document}
